%% sections/g-paygo-cost.tex - H. R. 9510 Bill v4.0 (full bill) - Appendix G
%% Genre: fiscal note (PAYGO and cost estimate). Source deliverable 07. Table 7.
\billsec{APPENDIX G.}{PAYGO AND COST ESTIMATE NOTE (NON-OPERATIVE).}\phantomsection\label{toc:aG}

\uscprov{1}{\textit{Independent research draft. This is an informal, pre-introduction
estimate, not an official Congressional Budget Office (CBO) cost estimate or score,
and is not endorsed by the FDA, HHS, the CBO, or any Member of Congress. Bill: H. R.
9510 (illustrative number; the Clerk assigns the real number at introduction), 119th
Congress, 2d Session. Refer to: House Committee on Energy and Commerce (Subcommittee
on Health).}}

\uscprov{1}{\textbf{Summary.} H. R. 9510 sets one sequencing rule for Physical AI
oncology clinical investigations: an automated verification, validation, and
uncertainty quantification process must clear robot-patient interaction code before
that code is generated or executed. The bill operates through the existing device
authorities of the Federal Food, Drug, and Cosmetic Act, adding a new section 515D
(21 U.S.C. 360e-5) together with conforming amendments. It creates no entitlement, no
direct spending program, no new tax, and no mandatory appropriation. Its budgetary
effects are discretionary: the Food and Drug Administration must issue final
regulations under section 515D and then fold a new sequencing check into existing
review and inspection workflows. Those costs are expected to be absorbable within
existing FDA device-program resources and the medical-device user-fee structure. On
that basis the bill is expected to be neutral under the Statutory Pay-As-You-Go Act of
2010, which reaches direct spending and revenue rather than discretionary agency
activity \cite{cbo-paygo, usc-360e4}.}

\uscprov{1}{\textbf{Basis of the estimate.} This note rests on the text of H. R. 9510
and on how its mechanism maps onto existing FDA authorities. The operative provision,
new section 515D, regulates the order of operations within investigations already
conducted under the device framework: classification under 21 U.S.C. 360c, premarket
review under 21 U.S.C. 360e and section 510(k) (21 U.S.C. 360(k)), and predetermined
change control plans under section 515C (21 U.S.C. 360e-4). The bill adds no program
office, grant authority, benefit, or collection that would alter Federal receipts. Its
cost is therefore the incremental agency effort to write one rulemaking and to embed a
sequencing check in existing review and inspection workflows. Because no formal CBO
estimate has issued, the assessments here are qualitative and conservative; informal
consultation with CBO can precede introduction, and a formal score typically follows
\cite{usc-360c, usc-360e, usc-360}.}

\uscprov{1}{\textbf{Direct spending and revenue (PAYGO).} Statutory PAYGO applies to
legislation affecting direct spending or revenue. H. R. 9510 does not establish or
expand any entitlement, mandatory grant, or other form of direct spending; its duties
are conditions on conduct, enforced through the existing prohibited-acts and
adulteration provisions of the Act (21 U.S.C. 331 and 21 U.S.C. 351) rather than
through new outlays. The bill imposes no tax, levies no new user fee by statute, and
changes no existing fee schedule, so it has no revenue effect. Any FDA implementation
funded through appropriated user-fee authority is treated as discretionary spending,
not direct spending, and is governed by appropriations rather than by the PAYGO
scorecard. The expected scorecard entry is therefore zero, and the bill is expected to
be PAYGO neutral \cite{usc-331, usc-351}.}

\uscprov{1}{\textbf{Discretionary cost (rulemaking and implementation).} The principal
cost driver is the direction in section 515D that the Secretary issue final
regulations not later than 365 days after enactment, with interim guidance permitted
in the meantime. Rulemaking of this scope, defining the binding standards, the gate
semantics (ACCEPT, BLOCK, ESCALATE), the verification-fraction and validation-agreement
floors, the documentation and attestation content, and the readiness gate, involves
notice-and-comment drafting, economic and paperwork analysis, and interagency review.
Once the rule takes effect, reviewers must read cleared verification records, confirm
readiness-gate determinations, and assess attestations during investigational and
premarket review, and inspectors must check sequencing compliance. These are
extensions of work the device program already performs on software, change control
plans, and quality systems under the Quality Management System Regulation (21 CFR part
820) and 21 CFR part 11, not a new function. The expected discretionary cost is modest
and absorbable within existing FDA device-program resources and the medical-device
user-fee structure, subject to confirmation in the formal CBO estimate and the
appropriations process \cite{fr-qmsr-2024, cfr-part820, cfr-part11}.}

\uscprov{1}{\textbf{Unfunded mandates.} Under the Unfunded Mandates Reform Act of
1995, the bill is not expected to impose an intergovernmental mandate above the
statutory threshold; it regulates device sponsors and investigators rather than
directing State, local, or tribal governments, and the savings clause of the Act (21
U.S.C. 360k) is preserved. The bill does impose a private-sector mandate on sponsors
of Physical AI oncology clinical investigations, who must run the verification process
before generating or executing robot-patient interaction code and must document and
attest to the result. The expected cost is the incremental documentation, review, and
attestation burden layered on validation activity that responsible sponsors already
perform, so the private-sector mandate cost is expected to fall below the threshold,
subject to CBO's formal determination \cite{usc-360k}.}

\uscprov{1}{\textbf{Earmark declaration (negative).} H. R. 9510 contains no
congressional earmarks, no limited tax benefits, and no limited tariff benefits as
those terms are defined in clause 9 of Rule XXI of the Rules of the House of
Representatives. This is a negative declaration: the bill makes no targeted
appropriation, revenue, or tariff provision benefiting a specific entity or limited
group. Its operative provisions are generally applicable regulatory requirements
administered by the FDA, and the bill carries no appropriation.}

\uscprov{1}{\textbf{Next step.} This note is a good-faith pre-introduction estimate
prepared to accompany H. R. 9510, not an official CBO product. The authoritative cost
estimate, covering any direct spending and revenue effects, the discretionary spending
projection over the standard budget window, and the formal mandate determination, is
produced by the Congressional Budget Office and typically follows introduction and
referral to House Energy and Commerce, with informal CBO consultation available
beforehand. Nothing here should be read to bind, anticipate, or substitute for CBO's
conclusions \cite{cbo-paygo, congress-howlawsmade}.}

\uscprov{1}{Table 7 sets out the line items of this pre-introduction estimate.}

{\footnotesize
\begin{xltabular}{\textwidth}{@{}L{4.6cm} Y@{}}
\hline\noalign{\vskip 0.04cm}
\textbf{Line item} & \textbf{Expected effect}\\
\hline\noalign{\vskip 0.04cm}
\endfirsthead
\hline\noalign{\vskip 0.04cm}
\textbf{Line item} & \textbf{Expected effect}\\
\hline\noalign{\vskip 0.04cm}
\endhead
\hline\noalign{\vskip 0.04cm}
\endlastfoot
Direct spending (mandatory) & None significant; no entitlement or mandatory program
created\\
\hline\noalign{\vskip 0.04cm}
Revenue & None; no new tax, fee mandate, or revenue change\\
\hline\noalign{\vskip 0.04cm}
Statutory PAYGO scorecard & Neutral; no significant direct spending or revenue
effect\\
\hline\noalign{\vskip 0.04cm}
Discretionary cost (rulemaking and implementation) & Modest; expected absorbable
within existing FDA device-program and user-fee resources\\
\hline\noalign{\vskip 0.04cm}
Intergovernmental mandate (UMRA) & None expected above the statutory threshold\\
\hline\noalign{\vskip 0.04cm}
Private-sector mandate (UMRA) & A mandate exists; expected documentation and review
costs are likely below the threshold\\
\hline\noalign{\vskip 0.04cm}
Earmarks, limited tax benefits, limited tariff benefits & None (negative
declaration)\\
\hline\noalign{\vskip 0.04cm}
Formal CBO score & Follows introduction and referral to House Energy and Commerce\\
\hline
\end{xltabular}}
\vspace{-0.8cm}
\figcaption{Table 7. PAYGO and cost line items for H. R. 9510. The bill directs
rulemaking only, creates no direct spending or revenue, and is expected to be PAYGO
neutral; the formal score is produced by the CBO after introduction.}
